\begin{definition}[Computational local constant]\label{mod:delta-finitedefinition}
Let $F$ be a nonarchimedean local field, let
$\chi:F^\times\to\mathbf C^\times$ have exact multiplicative conductor
$m\in\mathbf N$, and let $\psi:F\to\mathbf C^\times$ have exact additive
conductor $n\in\mathbf Z$ in Langlands's convention.  An
\emph{admissible denominator} is an element $\gamma\in F^\times$ satisfying
\[
 \operatorname{ord}_F(\gamma)=m+n.
\]
Equivalently,
$\gamma\mathcal O_F=\mathfrak p_F^{m+n}$.  Such denominators exist.  If
$\gamma$ and $\gamma'$ are admissible, their explicitly defined ratio
$a=\gamma'/\gamma$ lies in $U_F^0$ and satisfies $\gamma'=a\gamma$.

Put $Q_m=U_F^0/U_F^m$.  For $u\in U_F^0$ define
\[
 g_\gamma(u)=\psi(u/\gamma)\chi(u)^{-1}.
\]
If $u\equiv v\pmod {U_F^m}$, then $\chi(u/v)=1$ and
\[
 \frac{u-v}{\gamma}\in\mathfrak p_F^{-n},
\]
so $g_\gamma(u)=g_\gamma(v)$.  Thus the summand is first descended to a
function on $Q_m$ using this invariance proof; no representative of a
quotient class is selected.  Define the unnormalised finite Gauss sum
\[
 G_\gamma(\chi,\psi)
   =\sum_{[u]\in Q_m}\psi(u/\gamma)\chi(u)^{-1}.
\]
There is no factor involving $|Q_m|$, $q_F$, a sign, or a choice of Haar
measure.  The later integral--finite bridge supplies one positive common
coset measure, which disappears after applying the phase.

The sum is nonzero.  The formal proof is finite Fourier duality on
\[
 A_m=\mathcal O_F/\mathfrak p_F^m,
 \qquad
 T(x)=\sum_{[u]\in Q_m}\psi(xu/\gamma)\chi(u)^{-1}.
\]
When $m=0$, $Q_m$ is a singleton and its sole summand is nonzero.  Suppose
$m>0$ and $G_\gamma=T(1)=0$.  A unit frequency is reduced to $T(1)$ by
multiplicative reindexing.  A frequency in $\mathfrak p_F$ is killed by
reindexing with an element of $U_F^{m-1}$ on which $\chi$ is nontrivial;
exactness of the multiplicative conductor supplies that element.  Hence
$T(x)=0$ for every $x\in A_m$.  On the other hand, exactness of the additive
conductor makes the pairing
\[
 (x,y)\longmapsto\psi(xy/\gamma)
\]
on $A_m$ nondegenerate.  Finite additive-character orthogonality therefore
gives
\[
 \sum_{x\in A_m}T(x)\psi(-x/\gamma)=|A_m|\ne0,
\]
a contradiction.  In particular the proof does not assume nonvanishing
from a downstream integral or Lamprecht formula.

Multiplication by $a=\gamma'/\gamma$ permutes $Q_m$, and the descended
summand and the finite sum obey the exact change formulas
\[
 g_{\gamma'}(au)=\chi(a)^{-1}g_\gamma(u),
 \qquad
 G_{\gamma'}=\chi(a)^{-1}G_\gamma.
\]
The restriction of $\chi$ to $U_F^0$ factors through the finite group
$Q_m$, so $|\chi(a)|=1$; this proves the manuscript's unitary assertion
without a compactness or hidden-choice assumption.

With the total phase $\operatorname{ph}(0)=1$, define the computational
local constant by the manuscript normalization
\[
 \Delta_F^{\mathrm{fin}}(\chi,\psi;\gamma)
   =\chi(\gamma)\operatorname{ph}
      \bigl(G_\gamma(\chi,\psi)\bigr).
\]
The proved nonvanishing of $G_\gamma$ permits the nonzero branch of phase
multiplicativity.  Combining it with the two change factors above proves
unconditionally that this value is independent of the admissible
denominator $\gamma$.

\LeanFile{LanglandsFirstMainLemma/Delta/FiniteDefinition.lean}
\lean{LanglandsFirstMainLemma.AdmissibleGamma}
\lean{LanglandsFirstMainLemma.AdmissibleGamma.exists_admissible}
\lean{LanglandsFirstMainLemma.AdmissibleGamma.smul_lattice_zero}
\lean{LanglandsFirstMainLemma.AdmissibleGamma.ratio_mem_unitFiltration}
\lean{LanglandsFirstMainLemma.finiteGaussSummandRepresentative_eq_of_congruent}
\lean{LanglandsFirstMainLemma.finiteGaussSummand}
\lean{LanglandsFirstMainLemma.finiteGaussSummand_mk}
\lean{LanglandsFirstMainLemma.finiteGaussSum}
\lean{LanglandsFirstMainLemma.finiteGaussSum_ne_zero}
\lean{LanglandsFirstMainLemma.finiteGaussSum_admissibleGamma_change}
\lean{LanglandsFirstMainLemma.quasiChar_norm_eq_one_of_mem_unitFiltration_zero}
\lean{LanglandsFirstMainLemma.deltaFinite}
\lean{LanglandsFirstMainLemma.deltaFinite_gamma_independent}
\leanok
\SourceText{Definition of Delta and Lemma lem:gamma-independent.}
\uses{mod:basic-characterconductors,mod:localfield-finitequotients,mod:basic-phase}
\FormalizationContract The finite sum is unnormalised and is defined only
after its summand has been descended with an explicit congruence-invariance
proof.  Admissible denominators and quotient classes are explicit inputs;
no declaration chooses either one.  Nonvanishing is proved here because the
total phase cannot transport a zero sum through an admissible-denominator
change.
\end{definition}
